\documentclass[10pt,a4paper]{amsart}
\usepackage[margin=19mm]{geometry}
\usepackage{amsmath,amssymb,mathtools}
\usepackage[scr=rsfs,cal=euler]{mathalfa}
\usepackage{xfrac}
\usepackage[unicode,hidelinks,bookmarks=false]{hyperref}
\newcommand{\ie}{i.e.\ }
\newcommand{\cf}{cf.\ }
\newcommand{\resp}{respectively\ }
\newcommand{\Subsection}{\subsection}
\providecommand{\nobreakdash}{\nobreak\hbox{-}}
\setlength{\emergencystretch}{2em}
\multlinegap0pt
\usepackage{bm}
\usepackage{bbm}
\usepackage{dsfont}
\usepackage{xspace}
\usepackage{cancel}
\usepackage{mathtools}
\usepackage{amsthm}
\makeatletter
\@ifundefined{theorem}{\newtheorem{theorem}{Theorem}}{}
\@ifundefined{proposition}{\newtheorem{proposition}[theorem]{Proposition}}{}
\makeatother
\makeatletter
\expandafter\ifx\csname rhp\endcsname\relax
\newenvironment{rhp}[1][]{\refstepcounter{rhp}\par\medskip
   \noindent \textbf{Riemann-Hilbert problem~\therhp. #1} \rmfamily}{\medskip}
\fi
\expandafter\ifx\csname sproof\endcsname\relax
\newenvironment{sproof}{%
  \renewcommand{\proofname}{Sketch of the proof}\proof}{\endproof}
\fi
\makeatother
\title[A Tau function for $q$-Painlev\'e VI as a Fredholm determina]{A Tau function for $q$-Painlev\'e VI as a Fredholm determinant}
\author{Harini Desiraju, Pieter Roffelsen}
\date{Source: arXiv:2608.03345v1 (CC BY 4.0)}
\begin{document}
\maketitle

\noindent\emph{Source and licence.} A Tau function for $q$-Painlev\'e VI as a Fredholm determinant. Version arXiv:2608.03345v1 (CC BY 4.0), distributed under the Creative Commons Attribution 4.0 International licence, as recorded in the arXiv licence metadata for this exact version and shown on the abstract page https://arxiv.org/abs/2608.03345v1. Full rights notice: licenses/arxiv-2608.03345-CC-BY-4.0.txt.\par\smallskip

\noindent\textit{Editorial scope.} Counted unit: the Proposition whose source label is \texttt{prop:lax} in the article (source0.tex lines 621--655, source label \texttt{prop:lax}), delivered together with the article's own complete original proof of it (source0.tex lines 656--693, one \texttt{proof} environment of 38 lines). No mathematics is written, repaired or re-derived: statement and proof are the article's own text line for line. Required context retained uncounted: eq:det (lines 591--596). Every definition, display and label of those lines is retained unchanged. Omitted from this package and not redistributed: 1--590, 597--620, 694--2839. Nothing inside a retained excerpt is omitted or altered: every retained body is delivered exactly as the article's lines are written, so no omission receipt of any kind accompanies this package. Citations: the retained excerpts carry no citation command at all, so no bibliography entry of the article is transcribed and no reference is redistributed. Reference adaptation: none was needed and none was made; every reference inside the delivered unit resolves inside it, so no unresolved reference or citation is produced. Printed slips kept literal: no printed word, symbol, hypothesis or number of the counted statement or of its proof was altered. Layout and class: the article's own class is \texttt{amsart} with packages including amsfonts, tikz, tkz-graph, mathtools, geometry, accents, array, calc, bbold, mathrsfs, stmaryrd, pgf, tkz-euclide, caption; the wrapper is the lane's pinned \texttt{amsart} editorial wrapper at 10pt on A4, which restates the theorem-like environments the retained text needs and adds the article's own macro definitions that the retained text needs (none), with extra wrapper packages (bm, bbm, dsfont, xspace, cancel, mathtools). The delivered numbering is the unit's own. Assets: no retained span contains a figure environment, a graphics-inclusion command, a PDF-page inclusion, a table or a TikZ picture; no image file is copied into this package, the package declares no asset dependency and no image file is redistributed with it. Licence: version arXiv:2608.03345v1 of this article is distributed under the Creative Commons Attribution 4.0 International licence, as recorded in the arXiv licence metadata for this exact version and shown on the abstract page https://arxiv.org/abs/2608.03345v1. A full rights notice is delivered at licenses/arxiv-2608.03345-CC-BY-4.0.txt. Characters: every retained excerpt is pure ASCII, so the delivered engine, XeLaTeX with its default Latin Modern text font, reports no missing character.
\begin{equation}\label{eq:det}
\begin{aligned}
   |\Psi_\infty(z,t)|&=\left(q^{1+\theta_t}\frac{t}{z},q^{1-\theta_t}\frac{t}{z},q^{1+\theta_1}\frac{1}{z},q^{1-\theta_1}\frac{1}{z};q\right)_\infty,\\
   aq^{u(u-1)}t^2\times |\Psi_0(z,t)|&= \left(q^{+\theta_t}\frac{z}{t},q^{-\theta_t}\frac{z}{t},q^{+\theta_1}z,q^{-\theta_1}z;q\right)_\infty^{-1}.
\end{aligned}
\end{equation}

\begin{proposition}\label{prop:lax}
    The matrix functions
\begin{align*}
Y_0(z,t)&=z^{\log_q(-t)+\log_q(-1)}\Psi_0(z,t)z^{\theta_0\sigma_3},\\
Y_\infty(z,t)&=z^{\log_q(z/q)}\Psi_\infty(z,t) z^{-\theta_\infty\sigma_3},
\end{align*}
satisfy identical $q$-difference systems with respect to $z$ and $t$,
\begin{subequations}\label{eq:laxpair}
\begin{align}
   Y(qz,t)&=A(z,t)Y(z,t),\label{eq:laxpairA}\\
   Y(z,qt)&=B(z,t)Y(z,t). \label{eq:laxpairB}
\end{align}
\end{subequations}

The matrix function $A(z,t)$ is a degree two matrix polynomial in $z$,
\begin{equation}\label{eq:lin_sys}
   A(z,t)=A_0(t)+zA_1(t)+z^2 A_2
\end{equation}
with coefficients meromorphic in $t\in \mathcal{T}$ satisfying
\begin{equation}\label{eq:Amats}
A_2= q^{-\theta_\infty \sigma_3},\quad A_0(t)=H(t)t  q^{+\theta_0 \sigma_3}  H(t)^{-1},\quad H(t):=\Psi_0(0,t).
\end{equation}
The determinant of $A(z,t)$ is given by
\begin{equation}\label{eq:detA}
|A(z,t)|=(z-q^{+\theta_t}t)(z-q^{-\theta_t}t)(z-q^{+\theta_1})(z-q^{-\theta_1}).
\end{equation}
The function $B(z,t)$ in \eqref{eq:laxpairB} is a rational matrix function in $z$ of the form
\begin{equation*}
    B(z,t)=\frac{z^2I+zB_0(t)}{(z-q^{1+\theta_t}t)(z-q^{1-\theta_t}t)},\qquad B_0(t)=q^{2}t^{2}H(qt)H(t)^{-1},
\end{equation*}
where the matrix $B_0(t)$ is meromorphic in $t\in \mathcal{T}$, satisfying
\begin{equation*}
    |B_0(t)|=q^2 t^2,\qquad \operatorname{Tr} B_0(t)=-(q^{1+\theta_t}+q^{1-\theta_t})t.
\end{equation*}
\end{proposition}

\begin{proof}
Define the matrix function
\begin{align}
 A(z,t)&:=Y_\infty(qz,t)Y_\infty(z,t)^{-1}\nonumber\\
 &\,=  z^2 \Psi_\infty(qz,t)q^{-\theta_\infty\sigma_3}\Psi_\infty(z,t)^{-1} \label{eq:Amatrixdef}\\
 &\,= t\Psi_0(qz,t)q^{\theta_0\sigma_3}C(z,t)\Psi_\infty(z,t)^{-1}\nonumber\\
 &\,= t\Psi_0(qz,t)q^{\theta_0\sigma_3}\Psi_0(z,t)^{-1}\nonumber\\
 &\,=Y_0(qz,t)Y_0(z,t)^{-1}.\nonumber
\end{align}
It is a meromorphic function in $t\in \mathcal{T}$. Furthermore, it follows from the above expressions that $A(z,t)$ is analytic in $z\in \mathbb{C}$ and satisfies 
\begin{equation*}
    A(z,t)=z^2(q^{-\theta_\infty\sigma_3}+\mathcal{O}(z^{-1}))\qquad (z\rightarrow \infty),
\end{equation*}
and
\begin{equation*}
A(0,t)=\Psi_0(0,t)tq^{\theta_0\sigma_3}\Psi_0(0,t)^{-1}.
\end{equation*}
This shows that $A(z,t)$ is a degree two matrix polynomial in $z$. Finally the expression for the determinant of $A(z,t)$ follows from equations \eqref{eq:det} and the defining equations of $A(z,t)$,  \eqref{eq:Amatrixdef}.

Similarly, we define the $B$ matrix by
\begin{align*}
 B(z,t)&:=Y_\infty(z,qt)Y_\infty(z,t)^{-1}\\
 &\,=\Psi_\infty(z,qt)\Psi_\infty(z,t)^{-1}\\
 &\,=z\Psi_0(z,qt)C(z,t)\Psi_\infty(z,t)^{-1}\\
 &\,=z\Psi_0(z,qt)\Psi_0(z,t)^{-1}\\
 &\,=Y_0(z,qt)Y_0(z,t)^{-1}.
\end{align*}
It is a meromorphic function in $t\in \mathcal{T}$. Furthermore, it follows from the above expressions that it can only have poles in $z$ at $z=q^{1\pm \theta_t}t$, necessarily simple, and
\begin{align*}
 & &  B(z,t)&=I+\mathcal{O}(z^{-1}) &  &(z\rightarrow \infty), &\\
&  &  B(z,t)&=\mathcal{O}(z) & &(z\rightarrow 0). &
\end{align*}
Consequently, $B(z,t)$ is a rational matrix function in $z$ of the form indicated in the proposition, for some $B_0(t)$. The expressions for the determinant and trace of $B_0(t)$ follow from
\begin{equation*}
|zI+B_0(t)|= (z-q^{1+\theta_t}t)(z-q^{1-\theta_t}t),
\end{equation*}
which in turn follows from equations \eqref{eq:det}.
\end{proof}

\end{document}
